% Generated draft: requires mathematical and source review.
\documentclass[11pt]{article}
\usepackage[margin=25mm]{geometry}
\usepackage{fontspec}
\setmainfont{Latin Modern Roman}
\usepackage{amsmath,amssymb,mathtools,unicode-math}
\setmathfont{Latin Modern Math}
\usepackage{enumitem,needspace,xurl,hyperref}
\setlength{\parindent}{0pt}
\setlength{\parskip}{5pt}
\setlength{\emergencystretch}{3em}
\newcommand{\N}{\mathbb{N}}
\newcommand{\Z}{\mathbb{Z}}
\newcommand{\Q}{\mathbb{Q}}
\newcommand{\R}{\mathbb{R}}
\newcommand{\C}{\mathbb{C}}
\newcommand{\F}{\mathbb{F}}
\newcommand{\A}{\mathbb{A}}
\newcommand{\PP}{\mathbb{P}}
\newcommand{\E}{\mathbb{E}}
\newcommand{\eps}{\varepsilon}
\newcommand{\dd}{\,\mathrm{d}}
\begin{document}
\section*{UnsolvedMath ID: 30001025}
\textbf{Source problem code:} catalogue entry
\subsection{Definitions and mathematical statement}
For \(w\in[F,F]\), let \(\operatorname{cl}(w)\) be the least number of commutators whose product is \(w\), and define stable commutator length by
\[
 \operatorname{scl}(w)=\lim_{m\to\infty}\frac{\operatorname{cl}(w^m)}m.
\]
The question concerns the observed frequencies of rational values \(\operatorname{scl}(w)=p/q\) in lowest terms, with \(q>0\) and \(\gcd(p,q)=1\), among free-group words of a fixed length. The source does not specify the finite-rank free group, basis, word-reduction convention, probability measure, or restriction to the commutator subgroup. In a free group, words outside \([F,F]\) cannot be assigned a finite value by this commutator-length definition merely by taking a positive power.

For illustration of the needed sampling data, fix a finite basis of a free group and let \(\mathcal W_N\) be a chosen finite nonempty collection of length-\(N\) words with finite scl. Under uniform sampling one would define
\[
 \operatorname{freq}_N(p/q)=
 \frac{|\{w\in\mathcal W_N:\operatorname{scl}(w)=p/q\}|}{|\mathcal W_N|}.
\]
This is a conditional definition, not a selection of the source's missing word model. The proposed pattern is
\[
 \operatorname{freq}(p/q)\asymp q^{-\delta}
 \quad\text{for some exponent }\delta.
\]
In a rigorous asymptotic interpretation, \(\asymp\) would require positive upper and lower comparison constants on a specified range, with stated dependence on the rank, word length, numerator, and limiting regime. Those quantifiers are absent.

For any genuinely fixed finite word collection, only finitely many rational scl values occur, so there is no positive two-sided power law as \(q\to\infty\) across all denominators or all rational numbers. An eventual theorem would need, for example, a word-length limit and a specified growing denominator range, or a statement explicitly restricted to an empirical range. The source does not choose one. Also, frequencies of individual values \(p/q\) differ from the denominator marginal
\[
 \sum_{\substack{p\geq0\\\gcd(p,q)=1}}\operatorname{freq}_N(p/q),
\]
and the latter must not be substituted without source evidence.

The supplied assessment dated 2026-08-21 calls the proposed frequency power law an experimental conjectural pattern rather than a theorem. It also reports established scl consequences, including a negative answer to an unspecified Bavard question; it does not state that question, and those consequences are not a proof of the frequency law.
\subsection{Short English statement}
The source asks whether the observed frequency of a reduced rational scl value decays as a power of its denominator. A rigorous version needs a word distribution and an asymptotic regime: a fixed finite sample has only finitely many denominators, and individual-value frequencies are not denominator frequencies.
\subsection{Sources}
\begin{itemize}
\item Local UnsolvedMath registry, permanent ID \(30001025\), exact source code \texttt{catalogue entry}, proposed year 2008; no individual proposer is supplied.
\item \emph{Geometric Group Theory, Hyperbolic Dynamics and Symplectic Geometry (2009)}, Oberwolfach Reports, DOI 10.4171/owr/2008/42, pp. 2337--2406. \url{https://doi.org/10.4171/owr/2008/42}.
\item Calegari, \emph{scl}, \url{https://math.uchicago.edu/~dannyc/books/scl/scl.pdf}.
\item Calegari and Maher, ``statistics and compression of scl,'' \url{https://arxiv.org/abs/1008.4952}.
\item The supplied \texttt{corrected\_verified} note says the full statement was retained with mathematical delimiters. The partially solved label and assessment checked 2026-08-21 are supplied dated curation, not independently verified; established scl consequences and the experimental power-law claim must be reviewed separately.
\end{itemize}
\end{document}
